\documentclass[10pt]{article}
\usepackage[a4paper,margin=0.58in]{geometry}
\usepackage{amsmath,amssymb,amsthm}
\usepackage{enumitem}
\usepackage{microtype}

\setlist[itemize]{topsep=2pt,itemsep=2pt,leftmargin=*}
\setlist[enumerate]{topsep=2pt,itemsep=2pt,leftmargin=*}
\setlength{\parindent}{0pt}
\setlength{\parskip}{3pt}
\pagestyle{empty}

\title{\vspace{-2.5em}Preliminary note on Conjecture 3.2}
\author{Guangjian Zhang}
\date{}

\begin{document}
\maketitle
\vspace{-3em}
\small

\section*{Current understanding}

Let \(M\) be the initial phase space, let \(e\) denote the Euclidean metric, and let \(g_t=\phi_t^*e\) be the time-dependent pullback metric generated by a volume-preserving flow. On \(\mathbb M_0=[0,\tau]\times M\),
\[
G_{0,a}(t,x)=
\begin{pmatrix}a^{-2}&0\\0&g_t(x)\end{pmatrix},
\qquad
\Delta_{G_{0,a}}F(t,\cdot)=a^2\partial_{tt}F(t,\cdot)+\Delta_{g_t}F(t,\cdot).
\]

Let \(S_a,H_a\) be the Sobolev and Cheeger constants on \((\mathbb M_0,G_{0,a})\), and let \(s_D,h_D\) be the dynamic constants on \(M\). The known result is
\[
H_a=S_a\le s_D=h_D,
\]
with \(H_a,S_a\) nondecreasing in \(a\). Conjecture 3.2 asks whether
\[
\lim_{a\to\infty}S_a=s_D,\qquad
\lim_{a\to\infty}H_a=h_D,
\]
and whether minimizers converge to the time-copied dynamic minimizers. The inequality \(S_a\le s_D\) is supplied by testing \(S_a\) on time-independent functions \(F(t,x)=f(x)\). A possible way to view the remaining functional step is to prove a liminf lower bound, using compactness to show that near-minimizers become almost independent of time as \(a\to\infty\).

\section*{Possible first target}

I would first try to prove \(S_a\to s_D\) in a simple setting, for example on a compact boundaryless manifold or flat torus. Since \(S_a\le s_D\) is known, the key point is
\[
\liminf_{a\to\infty}S_a\ge s_D.
\]

After cancelling the common volume factor in \(dV_{G_{0,a}}=a^{-1}dt\,d\ell\), take a near-minimizing sequence \(U_a\). Since the quotient is invariant under adding constants and multiplying by nonzero scalars, one may replace \(U_a\) by the normalized representative
\[
F_a:=\frac{U_a-\alpha_a}{D(U_a)},
\qquad
D(U_a):=\inf_{\alpha\in\mathbb R}\int_0^\tau\int_M |U_a-\alpha|\,d\ell\,dt,
\]
where \(\alpha_a\) is a median or minimizing constant. Then \(F_a\) is still near-minimizing and satisfies
\[
\inf_{\alpha\in\mathbb R}\int_0^\tau\int_M |F_a-\alpha|\,d\ell\,dt=1.
\]
In the cancelled quotient, uniform boundedness of the numerator
\[
\int_0^\tau\int_M
\left(a^2|\partial_tF_a|^2+\|\nabla_{g_t}F_a(t,\cdot)\|_{g_t}^2\right)^{1/2}
d\ell\,dt
\le C
\]
gives
\[
\int_0^\tau\int_M|\partial_tF_a|\,d\ell\,dt\le C/a\to0.
\]
Thus a one-dimensional Poincare inequality in time should make \(F_a\) close in \(L^1([0,\tau]\times M)\) to \(\bar F_a(x)=\tau^{-1}\int_0^\tau F_a(t,x)\,dt\).

The step I am not yet sure about is the appropriate lower semicontinuity statement for the spatial total variation, roughly of the form
\[
\liminf_{a\to\infty}\int_0^\tau\int_M \|\nabla_{g_t}F_a(t,\cdot)\|_{g_t}\,d\ell\,dt
\ge
\int_0^\tau\int_M \|\nabla_{g_t}f\|_{g_t}\,d\ell\,dt,
\]
where \(F_a\to f(x)\) in \(L^1\) after passing to a subsequence. If the denominator also passes to the limit correctly, this would imply the desired liminf inequality. If this works, \(H_a\to h_D\) follows from \(H_a=S_a\) and \(s_D=h_D\). The convergence of minimizers could then be treated as a subsequent step after the functional convergence is understood.

\section*{Questions where guidance would be valuable}

\begin{itemize}
  \item Is Conjecture 3.2 still open, and is \(S_a\to s_D\) a reasonable first target?
  \item Are BV compactness and lower semicontinuity for weighted total variation the right tools, or is there a more natural route?
  \item Would the Dirichlet analogue, Conjecture 3.5, be a cleaner starting point, or does it introduce additional boundary difficulties?
\end{itemize}

If this first functional step is reasonable, I would try to write a more detailed note identifying exactly where compactness and lower semicontinuity are needed, before attempting to extend the argument toward the full conjecture.

\end{document}
